\section{পূর্ববর্তী বোর্ড প্রশ্নের ধারণাভিত্তিক অনুশীলন}

সংকলিত বোর্ড প্রশ্নের ধারণা অনুসারে নিচের অনুশীলনগুলো রূপান্তরিতভাবে সাজানো হয়েছে। উদ্দীপকের তথ্য থেকে কী, ডেটা টাইপ ও সম্পর্ক নির্ণয় করো এবং প্রতিটি সিদ্ধান্তের কারণ লেখো।

\subsection{কোন প্রশ্নে কোন ধারণা}
\begin{flushleft}
\small
\begin{tabularx}{0.98\linewidth}{|p{30mm}|X|}\hline
বোর্ড ও সাল & সংকলনে থাকা প্রশ্নের প্রধান চাহিদা\\ \hline
ঢাকা ২০১৯ & কম্পোজিট কী, এনক্রিপশন, ফিল্ডের টাইপ ও দুই টেবিলের সম্পর্ক\\ \hline
রাজশাহী ২০১৯ & RDBMS, সম্পর্কের শর্ত, নতুন ফিল্ড এবং জন্মতারিখে সর্টিং\\ \hline
কুমিল্লা ২০১৯ & DBMS, প্রাইমারি কী পরিবর্তন, টাইপ ও সম্পর্ক\\ \hline
যশোর ২০১৯ & ডেটা টাইপ, কর্পোরেট ডেটাবেজ ও সম্পর্কচিত্র\\ \hline
চট্টগ্রাম ২০১৯ & Ciphertext, প্রাইমারি/ফরেন কী, ভর্তি ও ফোন টেবিল\\ \hline
বরিশাল ২০১৯ & কর্পোরেট ডেটাবেজ, ছবি/অবজেক্ট, সম্পর্ক ও ইনডেক্সিং\\ \hline
সিলেট ২০১৯ & সংযোগ টেবিল, Lookup, একই রোলের জন্য সমন্বিত কী\\ \hline
দিনাজপুর ২০১৯ & ইনডেক্সের প্রয়োজন, Salary/Date/Photo-এর টাইপ ও টেবিল তৈরি\\ \hline
ঢাকা, যশোর, সিলেট, দিনাজপুর ২০১৮ & শিক্ষক ও রুটিন, SELECT-এর ব্যাখ্যা, কী ও সম্পর্ক\\ \hline
\end{tabularx}
\end{flushleft}
প্রতিটি উদ্দীপকে ব্যবহৃত পরিচিতি, ফিল্ডের মান ও প্রতিষ্ঠানের নিয়ম লক্ষ করো। একই ধরনের টেবিল থাকলেও নিয়ম বদলালে কী বা সম্পর্কের সিদ্ধান্ত বদলাতে পারে।

\subsection{বাকি সংকলিত প্রশ্নের প্রস্তুতি-মানচিত্র}
\begin{flushleft}
\small
\begin{tabularx}{0.98\linewidth}{|p{30mm}|X|}\hline
বোর্ড ও সাল & কী শিখে অনুশীলন করবে\\ \hline
রাজশাহী, কুমিল্লা, চট্টগ্রাম, বরিশাল ২০১৮ & নতুন ফিল্ড যোগ, কুয়েরি, গোপনীয়তা, দুটি টেবিলের সম্পর্ক\\ \hline
ঢাকা ২০১৭ & DBA-এর দায়িত্ব; তথ্য দিয়ে টেবিল নকশা; সুবিধা ও সীমা\\ \hline
রাজশাহী ২০১৭ & SQL-এর ভূমিকা; শর্তে বাছাই; সংযোগকারী কী যোগ; ইনডেক্স\\ \hline
যশোর ২০১৭ & ইনডেক্স হালনাগাদ, টাইপ নির্বাচন, সম্পর্কের যুক্তি\\ \hline
কুমিল্লা ২০১৭ & একই নামের ব্যক্তিকে অনন্য কী দিয়ে খোঁজা; সম্পর্ক\\ \hline
চট্টগ্রাম ২০১৭ & রেকর্ড, এনক্রিপশন, কী বাদ দিলে সমস্যা, সম্পর্কের শর্ত\\ \hline
সিলেট ২০১৭ & SQL, ইনডেক্সের ব্যবহার, জন্মতারিখের টাইপ, সম্পর্ক\\ \hline
বরিশাল ২০১৭ & SQL, ইনডেক্স, সংখ্যামান ও গ্রেডের টাইপ, সম্পর্কের সুবিধা\\ \hline
দিনাজপুর ২০১৭ & টেবিল তৈরি, সুবিধা-সীমা; অনন্য ফিল্ড, DBMS-এর কাজ\\ \hline
\end{tabularx}
\end{flushleft}
সারণিতে ২০১৭--২০১৯ সালের সংকলিত প্রশ্নে ব্যবহৃত ধারণাগুলো একত্রে দেখানো হয়েছে। একই ধারণা ভিন্ন উদ্দীপকে কীভাবে প্রয়োগ করতে হয়, তা পরের অনুশীলনগুলোতে লক্ষ করো।

\textbf{সিলেট ২০১৯-এর ধারণা প্রয়োগ:} একই Roll XI ও XII-তে থাকলে শুধু Roll নয়, নির্দিষ্ট শর্তে ClassName ও Roll-এর সমন্বয় ব্যবহার করো। Lookup তালিকা এন্ট্রি সহজ করে; এটি নিজে একটি নতুন মৌলিক ডেটা টাইপ নয়।

\textbf{দিনাজপুর ২০১৯-এর ধারণা প্রয়োগ:} EmployeeID, Salary, JoiningDate, Photo নিয়ে নকশা করো। Access-এ Salary--Currency, JoiningDate--Date/Time, Photo--Attachment/OLE Object হতে পারে। SQLite-এ আলাদা সংরক্ষণরীতি লাগে; যোগদানের তারিখের জন্য নির্দিষ্ট বিন্যাসের TEXT এবং ছবির বাইনারি ডেটায় BLOB ব্যবহার করা যায়। টাইপের সঙ্গে সফটওয়্যারের নামও লেখো।

\textbf{রাজশাহী ২০১৭-এর ধারণা প্রয়োগ:} দুই টেবিলে শুধু একই সারি-অবস্থান দেখে ব্যক্তির তথ্য জোড়া দেবে না। দ্বিতীয় টেবিলে সঠিক ব্যক্তির ID বসিয়ে সম্পর্ক তৈরি করো; ইচ্ছামতো পরিচিতি বানিয়ে দিলে ভুল ব্যক্তির বেতন যুক্ত হতে পারে।

\textbf{চট্টগ্রাম ২০১৭-এর ধারণা প্রয়োগ:} Name পুনরাবৃত্ত থাকলে পরিচিতির কী বাদ দিলে নির্দিষ্ট রেকর্ড নির্বাচন ও অন্য টেবিলের সঙ্গে সংযোগে দ্ব্যর্থতা তৈরি হয়। শুধু নামকে কী ধরে সমস্যার সমাধান হবে না।

\subsection{অনুশীলন ১: দুই টেবিলে একজন শিক্ষার্থী}
\textbf{ধারণাসূত্র: ঢাকা ২০১৯—রূপান্তরিত অনুশীলন।}
একটি কলেজে Student(Roll, Name, BirthDate, ClassName) এবং Profile(Roll, GuardianName, GPA) আছে। এখানে প্রত্যেক শিক্ষার্থীর একটি Profile এবং দুই টেবিলেই Roll অনন্য।
\begin{multicols}{2}
\small
\begin{enumerate}
\item[ক.] একাধিক ফিল্ড মিলিয়ে কী তৈরি করার প্রয়োজন কখন হয়?
\item[খ.] ব্যক্তিগত তথ্য গোপন রাখতে এনক্রিপশন কীভাবে কাজ করে?
\item[গ.] Name, BirthDate ও GPA-এর টাইপ নির্বাচন করে যুক্তি দাও।
\item[ঘ.] Student ও Profile-এর সম্পর্কের ধরন, কী এবং বাস্তবায়নের শর্ত ব্যাখ্যা করো।
\end{enumerate}
\small

\end{multicols}


\subsection{অনুশীলন ২: জন্মতারিখ ও সর্টিং}
\textbf{ধারণাসূত্র: রাজশাহী ২০১৯—রূপান্তরিত অনুশীলন।}
\begin{flushleft}
\small
\begin{tabular}{|l|l|l|}\hline
CustomerID & District & BirthDate\\ \hline
201 & Dhaka & 2007-10-15\\\hline
202 & Khulna & 2006-09-05\\\hline
203 & Bhola & 2008-12-20\\ \hline
\end{tabular}
\end{flushleft}
\begin{enumerate}
\item[ক.] RDBMS-এর পরিচয় দাও।
\item[খ.] দুটি ফিল্ডের মধ্যে সম্পর্ক স্থাপনে টাইপের সামঞ্জস্য কেন প্রয়োজন?
\item[গ.] Access-এ Booking টেবিলে BookingDate নামে Date/Time ফিল্ড যোগ করার ধাপ লেখো।
\item[ঘ.] জন্মতারিখ ঊর্ধ্বক্রমে সাজাও এবং বয়সে কে বড় তা ব্যাখ্যা করো।
\end{enumerate}

\begin{multicols}{2}
\small
\begin{enumerate}
\item[ক.] সংকেতিত বার্তা বলতে কী বোঝায়?
\item[খ.] প্রাইমারি কী ও ফরেন কী-এর কাজ আলাদা কেন?
\item[গ.] AdmissionDate, Fee ও Mobile-এর Access টাইপ নির্বাচন করো।
\item[ঘ.] কী ও সম্পর্কচিত্র দেখিয়ে নকশাটি বিশ্লেষণ করো।
\end{enumerate}
\small

\end{multicols}


\subsection{অনুশীলন ৪: শিক্ষক ও রুটিন}
\textbf{ধারণাসূত্র: ঢাকা, যশোর, সিলেট ও দিনাজপুর ২০১৮—রূপান্তরিত অনুশীলন।}
Teacher(TID, Name, SubjectName) এবং Routine(RoutineID, TID, GroupName, ClassTime) আছে। একই শিক্ষক বিভিন্ন সময়ে একাধিক ক্লাস নেন।
\textbf{প্রশ্ন:} Teacher থেকে Routine-এর সম্পর্ক দেখাও। পৃথক Students টেবিলে শিক্ষার্থীদের Roll ও Name আছে ধরে, শুধু ওই দুটি ফিল্ড নির্বাচন করার কুয়েরি কী করবে, ব্যাখ্যা করো।



\subsection{অনুশীলন ৫: কী পরিবর্তন ও সংশ্লিষ্ট তথ্য}
\textbf{ধারণাসূত্র: কুমিল্লা ২০১৯—রূপান্তরিত অনুশীলন।}
Student(Roll, Name, BirthDate) ও Result(Roll, ICTMark) টেবিলে প্রতিটি রোল একবার আছে। Result.Roll হলো Student.Roll-এর ফরেন কী। একজন কর্মী Student-এর Roll 301-কে 401 করতে চান, কিন্তু Result-এ 301 রয়েছে।
\begin{multicols}{2}
\small
\begin{enumerate}
\item[ক.] DBMS কী?
\item[খ.] প্রাইমারি কী পরিবর্তনে সংশ্লিষ্ট তথ্য বিবেচনা করতে হয় কেন?
\item[গ.] Name, BirthDate ও ICTMark-এর Access টাইপ নির্বাচন করো।
\item[ঘ.] সম্পর্ক অক্ষুণ্ণ রেখে পরিবর্তনের শর্ত এবং ভুল পরিবর্তনের পরিণতি ব্যাখ্যা করো।
\end{enumerate}
\small

\end{multicols}


\subsection{অনুশীলন ৬: ফল, ঠিকানা ও ER উপস্থাপন}
\textbf{ধারণাসূত্র: যশোর ২০১৯—রূপান্তরিত অনুশীলন।}
একটি কলেজে Academic(StudentID, Name, GPA, Remark) এবং Contact(StudentID, District) আছে। প্রতি শিক্ষার্থীর প্রতিটি টেবিলে সর্বোচ্চ একটি রেকর্ড। বৃহত্তর ব্যবস্থায় ভর্তি, হিসাব ও প্রশাসন বিভাগ এই তথ্য প্রয়োজনমতো ব্যবহার করে।
\begin{enumerate}
\item[ক.] ডেটা টাইপ কী?
\item[খ.] প্রতিষ্ঠানের বিভাগগুলোর জন্য সমন্বিত ডেটাবেজ প্রয়োজন কেন?
\item[গ.] Academic-এর ফিল্ডগুলোর টাইপ ও প্রয়োজনীয় বিধি নির্ধারণ করো।
\item[ঘ.] দুটি টেবিলের সম্পর্ক ER ধারণা ব্যবহার করে দেখাও।
\end{enumerate}
\subsection{অনুশীলন ৭: বিক্রয়, সম্পর্ক ও ইনডেক্স}
\textbf{ধারণাসূত্র: বরিশাল ২০১৯—রূপান্তরিত অনুশীলন।}
Product(ProductID, ProductName, UnitPrice) এবং Sale(SaleID, ProductID, Quantity) আছে। একটি পণ্যের বহু বিক্রয় হয়। ব্যবস্থাপক ProductID দিয়ে দ্রুত বিক্রয় খুঁজতে চান এবং দাম অনুযায়ী পণ্যের তালিকাও চান।
\begin{multicols}{2}
\small
\begin{enumerate}
\item[ক.] কর্পোরেট ডেটাবেজ কী?
\item[খ.] Access-এ ছবি/অবজেক্ট রাখার সঙ্গে Text নাম রাখার পার্থক্য কী?
\item[গ.] Product ও Sale-এর কী ও সম্পর্ক নির্ধারণ করো।
\item[ঘ.] দ্রুত খোঁজা ও দাম অনুযায়ী তালিকার জন্য ইনডেক্সিং ও সর্টিংয়ের ভূমিকা বিশ্লেষণ করো।
\end{enumerate}
\small

\end{multicols}


\subsection{অনুশীলন ৮: একই রোল ও যৌগিক কী}
\textbf{ধারণাসূত্র: সিলেট ২০১৯—রূপান্তরিত অনুশীলন।}
এক শিক্ষাবর্ষে Learner(Name, Roll, ClassName, District)-এ XI ও XII উভয় শ্রেণিতে Roll 1 ও 2 আছে। প্রতিটি শ্রেণির একটিমাত্র শাখা এবং শ্রেণির ভেতরে রোল অনন্য। District নির্দিষ্ট তালিকা থেকে নেওয়া হবে।
\begin{multicols}{2}
\small
\begin{enumerate}
\item[ক.] জাংশন টেবিল কী?
\item[খ.] Lookup তালিকা এন্ট্রি সহজ করলেও সব ভুল দূর করে না কেন?
\item[গ.] ফিল্ডগুলোর যুক্তিযুক্ত Access টাইপ নির্ধারণ করো।
\item[ঘ.] বর্তমান নিয়মে প্রাইমারি কী নির্বাচন এবং তার সীমা বিশ্লেষণ করো।
\end{enumerate}
\small

\end{multicols}


\subsection{অনুশীলন ৯: কর্মচারীর টেবিল তৈরি}
\textbf{ধারণাসূত্র: দিনাজপুর ২০১৯—রূপান্তরিত অনুশীলন।}
একটি প্রতিষ্ঠানে EmployeeID, Salary, JoiningDate ও Photo সংরক্ষণ করতে হবে। প্রশ্নে SQLite ব্যবহার করতে বলা হয়েছে; Salary টাকায় দুই দশমিক পর্যন্ত প্রয়োজন।
\begin{enumerate}
\item[ক.] ডেটাবেজ কী?
\item[খ.] নির্দিষ্ট পরিচিতি দিয়ে নিয়মিত খোঁজায় ইনডেক্স কেন সহায়ক?
\item[গ.] ফিল্ডগুলোর সংরক্ষণরীতি নির্বাচন করো এবং Access-এর সঙ্গে পার্থক্য লেখো।
\item[ঘ.] গ্রহণযোগ্য SQLite CREATE TABLE নির্দেশ লেখো।
\end{enumerate}
\subsection{বিগত MCQ-এর ধারণা: এইচএসসি ২০১৮}
\textbf{এইচএসসি ২০১৮-এর সংকলিত MCQ-এর ধারণাভিত্তিক অনুশীলন।} প্রশ্নগুলোর উত্তর দেওয়ার সময় কুয়েরি, সর্টিং ও নিরাপত্তার পৃথক উদ্দেশ্য বিবেচনা করো।
\begin{enumerate}
\item Access-এ Text ও Currency ফিল্ডের মান অনুযায়ী রেকর্ড সাজানো যায় কি?\\
\item কলেজের ফল থেকে নির্দিষ্ট শর্তের শিক্ষার্থীর নাম আলাদা করতে কোন কাজটি দরকার?\\
\item DBMS দিয়ে তথ্যের বিন্যাস ও উপস্থাপনা করা যায়—তাই কি নিরাপত্তা স্বয়ংক্রিয়ভাবে নিশ্চিত?\\
\item গোপনীয়তার উদ্দেশ্যে মূল বার্তাকে সংকেতিত করে কোন প্রক্রিয়া?\\
\end{enumerate}
